\chapter{Introducción}

Un párrafo con la descripción del contenido, lo que espera encontrar el lector.

\section{Antecedentes breves del problema}
Aquí se presenta una descripción breve del problema abordado por el Proyecto Integrador. Se recomienda contextualizar la temática y mencionar desarrollos previos relevantes.

\section{Relevancia del trabajo}
En esta sección se justifica la importancia del trabajo realizado, su aporte y el interés de su desarrollo dentro del área correspondiente.

\section{Motivación para la elección del tema}
Aquí se exponen las razones académicas, técnicas o personales que motivaron la elección del tema.

\section{Formulación del problema}
Se plantea claramente el problema de ingeniería que se pretende resolver, incluyendo restricciones y requerimientos principales.

\section{Objetivo general y objetivos específicos}
\subsection{Objetivo general}
Redactar aquí el objetivo general.

\subsection{Objetivos específicos}
\begin{itemize}
    \item Objetivo específico 1.
    \item Objetivo específico 2.
    \item Objetivo específico 3.
\end{itemize}

\section{Metodología utilizada para lograr los objetivos propuestos}
Describir brevemente la metodología general empleada para el desarrollo del trabajo.

\section{Orientación al lector de la organización del texto}
Indicar cómo está organizado el documento y qué encontrará el lector en cada parte.